\subsubsection{2.3. Mạng nơ-ron đồ thị (GNN)}
\indent\indent Mạng nơ-ron đồ thị là một lớp mô hình học sâu chuyên biệt được thiết kế để khai thác dữ liệu trong không gian phi Euclid. Theo đánh giá tổng quan của Zhang và cộng sự \cite{zhang2019graph}, mục tiêu của GNN là học các vector biểu diễn nút (\textit{node embeddings}) cấp thấp sao cho bảo tồn được cả đặc trưng nội tại của nút và cấu trúc liên kết phức tạp của mạng lưới. Trong bài toán này, GNN đóng vai trò là "bộ não" trích xuất các mẫu hành vi của các cụm Sybil từ đồ thị đa tầng.\vspace{0.3cm}

\begin{figure}[H]
    \centering
    \includegraphics[width=0.8\linewidth]{images/part-02/chap-01/GNN.png}
    \vspace{0.5cm}
    \caption{Minh hoạ kiến trúc cơ bản của GNN}
    \label{fig:GNN}
\end{figure}

Cơ chế hoạt động cốt lõi của GNN là truyền thông điệp (\textit{Message Passing}). Tại mỗi tầng $l$, đặc trưng của nút $v$ được cập nhật thông qua việc tổng hợp thông tin từ tập hàng xóm $N(v)$:\vspace{-0.1cm}
\[
    h_v^{(l)} = \text{UPDATE}^{(l)} \left( h_v^{(l-1)}, \text{AGGREGATE}^{(l)} \left( \{ h_u^{(l-1)} | u \in N(v) \} \right) \right)
\]
Trong đó $h_v^{(l)}$ là trạng thái ẩn của nút tại tầng $l$.

\begin{figure}[H]
    \centering
    \includegraphics[width=0.8\linewidth]{images/part-02/chap-01/GNN_MessagePassing.png}
    \vspace{0.5cm}
    \caption{Cơ chế truyền tin cơ bản trong mạng GNN}
    \label{fig:GNN_MP}
\end{figure}

\textbf{2.3.1. Các kiến trúc mạng nơ-ron đồ thị thực nghiệm}\vspace{0.2cm}

Với đặc thù phức tạp của mạng lưới SocialFi, đề tài khôn    g giới hạn ở một kiến trúc duy nhất mà thiết lập một hệ thống thử nghiệm toàn diện. Năm kiến trúc GNN tiên tiến được triển khai và đánh giá để tìm ra mô hình nhận diện tài khoản giả mạo tối ưu nhất:

\begin{itemize}[label={--}, itemsep=5pt]
    \item \textbf{Graph Attention Networks (GATv2):} Đóng vai trò là kiến trúc chủ đạo của hệ thống. GAT nguyên bản \cite{velickovic2017graph} tạo đột phá bằng cơ chế tự chú ý, nhưng bộc lộ giới hạn về "chú ý tĩnh", khiến mức độ quan trọng của láng giềng bị đánh giá toàn cục. Khắc phục rào cản này, Brody và cộng sự \cite{brody2021attentive} đã đề xuất GATv2 với cơ chế "chú ý động". Nhờ cơ chế chú ý động của GATv2, mô hình trở nên linh hoạt hơn trong việc phân bổ trọng số, từ đó bước đầu học được cách dồn sự chú ý vào các liên kết mang tín hiệu quyết định. Bản chất toán học của kiến trúc GATv2 nguyên bản được thể hiện qua các bước sau:
    \begin{itemize}[label={+}, itemsep=2pt]
        \item \textbf{Hệ số chú ý chưa chuẩn hóa} $e_{ij}$ được tính bằng cách ghép nối vector đặc trưng của hai nút, biến đổi tuyến tính và áp dụng hàm phi tuyến trước khi nhân với vector chú ý:
        \[
            e_{ij} = \vec{a}^T \text{LeakyReLU}\left( \mathbf{W} [h_i \| h_j] \right)
        \]
        Trong đó: $h_i, h_j$ là vector đặc trưng ẩn của nút gốc $i$ và láng giềng $j$; $\mathbf{W}$ là ma trận trọng số học được; $\vec{a}$ là vector chú ý.
        
        \item \textbf{Trọng số chú ý} $\alpha_{ij}$ được chuẩn hóa bằng hàm softmax để có thể so sánh giữa các láng giềng:
        \[
            \alpha_{ij} = \frac{\exp(e_{ij})}{\sum_{k \in N_i} \exp(e_{ik})}
        \]
        
        \item \textbf{Vector đầu ra} được tổng hợp bằng cách tính tổng có trọng số và đưa qua hàm phi tuyến, trong nghiên cứu này là hàm ELU:
        \[
            \vec{h}'_i = \text{ELU} \left( \sum_{j \in N_i} \alpha_{ij} \mathbf{W}\vec{h}_j \right)
        \]
        
        \item Mô hình tích hợp \textbf{\textit{Multi-head Attention}} tương tự Transformer \cite{vaswani2017attention} để ổn định quá trình học, ghép nối ở tầng ẩn và lấy trung bình ở tầng đầu ra.
    \end{itemize}

    \item \textbf{Graph Convolutional Networks (GCN):} Dựa trên công trình nền tảng của Kipf và Welling \cite{kipf2016variational}, GCN sử dụng phép xấp xỉ bậc nhất của tích chập phổ đồ thị để lan truyền và tổng hợp thông tin. Trong khung thử nghiệm của đề tài, mô hình này được triển khai để đóng vai trò là thước đo đối chứng cơ sở. Khác biệt cốt lõi với cơ chế chú ý động của GAT là GCN sử dụng cơ chế tổng hợp mang đặc trưng của một bộ lọc tĩnh. Tính chất này được thể hiện qua công thức lan truyền đa tầng ở dạng ma trận:
    \[
        H^{(l+1)} = \sigma\left( \tilde{D}^{-1/2} \tilde{A} \tilde{D}^{-1/2} H^{(l)} \mathbf{W}^{(l)} \right)
    \]
    Trong đó, $H^{(l)}$ là ma trận đặc trưng của các nút tại tầng $l$; $\mathbf{W}^{(l)}$ là ma trận trọng số học được; $\tilde{A}$ là ma trận kề đã thêm vòng lặp tự thân (\textit{self-loops}) và $\tilde{D}$ là ma trận bậc đường chéo tương ứng.\vspace{0.3cm}

    Bản chất "tĩnh" và sự hạn chế của GCN nguyên bản được bộc lộ rõ nhất khi xét quá trình lan truyền thông điệp từ láng giềng $j$ đến nút $i$ ở cấp độ từng nút:
    \[
        h_i^{(l+1)} = \sigma \left( \sum_{j \in \mathcal{N}(i) \cup \{i\}} \frac{1}{\sqrt{\tilde{d}_i \tilde{d}_j}} \mathbf{W}^{(l)} h_j^{(l)} \right)
    \]

    \item \textbf{Graph Isomorphism Network with Edge features (GINE):} Đóng vai trò là kiến trúc đại diện cho trường phái phân tích cấu trúc tô-pô chuyên sâu. Nguyên bản của kiến trúc này là mạng GIN (\textit{Graph Isomorphism Network}) được đề xuất bởi Xu và cộng sự \cite{xu2018powerful}. Điểm cốt lõi giúp nó vượt trội hơn GCN nằm ở cơ chế tổng hợp láng giềng mang tính đơn ánh trên các đa tập hợp. Thay vì sử dụng phép lấy trung bình, kiến trúc này sử dụng phép tổng kết hợp với mạng nơ-ron đa tầng (MLP). Tuy nhiên, GIN nguyên bản bộc lộ giới hạn lớn khi không có khả năng xử lý các đặc trưng cạnh. Từ đó biến thể GINE (\textit{GIN with Edges features}) được giới thiệu thông qua nghiên cứu của Hu và cộng sự \cite{hu2019strategies}. Công thức truyền tin được tùy biến để nhúng trực tiếp trọng số cạnh $e_{uv}$ vào thông điệp của láng giềng trước khi thực hiện phép gộp:
    \[
        h_v^{(k)} = \text{MLP}^{(k)} \left( (1 + \epsilon^{(k)}) h_v^{(k-1)} + \sum_{u \in \mathcal{N}(v)} \text{ReLU} \left( h_u^{(k-1)} + e_{uv} \right) \right)
    \]
    Trong đó, $e_{uv}$ chính là trọng số vô hướng $w_{ij}$ đại diện cho chi phí hành vi. Việc thực hiện thực nghiệm trên GINE không chỉ cho phép nhận diện được các cấu trúc hình học dị thường của các cụm Sybil (như mạng lưới hình sao), mà còn phân biệt được bản chất của các liên kết đó.

    \item \textbf{Principal Neighbourhood Aggregation (PNA):} Để khắc phục tình trạng chênh lệch bậc cực lớn, PNA được đề xuất bởi Corso và cộng sự \cite{corso2020principal}, từ bỏ cơ chế tổng hợp truyền thống. Thay vào đó, lõi toán học của mô hình tối ưu hóa việc trích xuất phân phối thông tin thông qua công thức tổng hợp láng giềng:
    $$
        h_i^{(l+1)} = \text{MLP} \left( h_i^{(l)}, \bigoplus_{j \in \mathcal{N}(i)} h_j^{(l)} \right)
    $$
    Trong đó, toán tử tổng hợp $\oplus$ là sự kết hợp của 4 bộ thống kê (Mean, Max, Min, Std) và 3 bộ chia tỷ lệ theo bậc dựa trên hàm logarit:
    $$
        \oplus = \begin{bmatrix} I \\ S(d, \alpha=1) \\ S(d, \alpha=-1) \end{bmatrix} \otimes \begin{bmatrix} \mu \\ \sigma \\ \max \\ \min \end{bmatrix}
    $$
    Tích tensor này cung cấp một không gian biểu diễn đa diện gồm 12 chiều cho mỗi thông điệp lan truyền. Nhờ các bộ chia tỷ lệ, mô hình có khả năng tự động điều chỉnh độ lớn của tín hiệu tùy thuộc vào bậc kết nối của từng nút. Đặc tính này mang lại lợi thế lớn, không bị đánh lừa bởi các cụm Sybil có số lượng liên kết khổng lồ với các hành vi rập khuôn, nghèo nàn về mặt phân phối (độ lệch chuẩn $\sigma$ cực thấp).
    
    \item \textbf{Graph Transformer Networks:} Được chuyển thể từ cơ chế tự chú ý (\textit{Self-Attention}) mang tính bước ngoặt của mô hình Transformer thuần túy trong lĩnh vực xử lý ngôn ngữ tự nhiên \cite{vaswani2017attention}, kiến trúc này được tùy biến cho đồ thị thông qua mạng truyền tin hợp nhất UniMP \cite{shi2020masked}. Khác biệt hoàn toàn với các bộ lọc cấu trúc (như GCN hay GINE) hay bộ thống kê (PNA), Graph Transformer thực hiện tổng hợp láng giềng dựa trên cơ chế chú ý tích vô hướng có tỷ lệ:
    \[
        h_i^{(l+1)} = \sum_{j \in \mathcal{N}(i)} \alpha_{ij} \mathbf{W}_v h_j^{(l)} \quad \text{với} \quad \alpha_{ij} = \text{softmax} \left( \frac{(\mathbf{W}_q h_i)^T (\mathbf{W}_k h_j)}{\sqrt{d}} \right)
    \]
    Trong đó, $\mathbf{W}_q, \mathbf{W}_k, \mathbf{W}_v$ lần lượt là các ma trận phóng chiếu Truy vấn (\textit{Query}), Khóa (\textit{Key}) và Giá trị (\textit{Value}). Phép nhân vô hướng tử số đo lường trực tiếp độ tương đồng sâu sắc giữa nút trung tâm và từng láng giềng. Cơ chế gán trọng số động ($\alpha_{ij}$) này đóng vai trò như một bộ lọc. Graph Transformer có thể dễ dàng cắt đứt sự chú ý ($\alpha_{ij} \approx 0$) đối với các liên kết mang tính tấng công khi phát hiện hồ sơ của chúng hoàn toàn "lệch pha" so với cộng đồng người dùng thực tế.
\end{itemize}
\vspace{0.3cm}

%=============================================================================

\textbf{2.3.2. Tối ưu hóa tham số cho đồ thị đa quan hệ}\vspace{0.2cm}

Với kiến trúc đồ thị đa quan hệ có trọng số, các thực thể tương tác thông qua 18 loại quan hệ khác nhau. Nếu biểu diễn các cạnh này theo phương pháp mã hóa One-hot 18 chiều và đưa trực tiếp vào mạng GNN, hệ thống sẽ gặp phải tình trạng bùng nổ tham số. Hiện tượng quá khớp trên tập dữ liệu thưa thớt. Để giải quyết triệt để rào cản này, đề tài thực hiện chuỗi tối ưu hóa tham số từ mức độ dữ liệu đến lõi kiến trúc:

\begin{itemize}[label={--}, itemsep=5pt]
        \item \textbf{Tối ưu kiến trúc GATv2 (Edge-augmented GATv2):} Để kiến trúc này tận dụng tối đa tham số $w_{ij}$ vừa được nén, đề tài thiết lập cơ chế \textit{dung hợp sớm trọng số cạnh}. Trọng số cạnh $w_{ij}$ được tích hợp trực tiếp vào quá trình tính toán hệ số chú ý thông qua phép cộng tuyến tính trước hàm kích hoạt phi tuyến:
        \[
            e_{ij} = \vec{a}^T \text{LeakyReLU}\left( \mathbf{W}_n [h_i \| h_j] + \mathbf{W}_e w_{ij} \right)
        \]
        Trong đó, $\mathbf{W}_e$ là vector tham số học được đóng vai trò phóng chiếu trọng số vô hướng $w_{ij}$ vào cùng không gian ẩn với đặc trưng nút. Việc tối ưu hóa bằng cách đặt tham số cạnh vào bên trong hàm LeakyReLU ép mạng nơ-ron phải dẫn dắt sự chú ý dựa trên tương tác chéo giữa bản chất thực thể và chi phí hành vi ($w_{ij}$).

        \item \textbf{Tinh chỉnh các kiến trúc cơ sở:} Việc đồng bộ hóa không gian tham số đầu vào không chỉ dừng lại ở bước tiền xử lý mà là sự can thiệp trực tiếp vào lõi toán học của toàn bộ các mô hình đối chứng. Dựa trên đặc thù dữ liệu của đồ thị đa quan hệ có trọng số hiện tại, mỗi kiến trúc buộc phải thay đổi cơ chế tổng hợp láng giềng để tiếp nhận trọng số hành vi $w_{ij}$:

        \begin{itemize}[label={+}, itemsep=4pt]

            \item \textbf{Graph Convolutional Networks (GCN):} Để xóa bỏ cơ chế cào bằng theo cấu trúc, hệ thống can thiệp vào luồng truyền tin bằng cách truyền trực tiếp tensor \texttt{edge\_weight} 1D. Ma trận kề nhị phân nguyên thủy bị vô hiệu hóa, hệ số chuẩn hóa tĩnh được thay thế bằng trọng số hành vi, dẫn đến công thức cập nhật mới:
            \[
                h_i^{(l+1)} = \sigma \left( \sum_{j \in \mathcal{N}(i) \cup \{i\}} \frac{w_{ij}}{\sqrt{\tilde{d}_i \tilde{d}_j}} \mathbf{W}^{(l)} h_j^{(l)} \right)
            \]
            Trong đó, bậc $\tilde{d}_i$ được tính dựa trên tổng trọng số các liên kết kề. Phép toán này tước bỏ lợi thế của các nút siêu đặc (vốn dĩ có số lượng kết nối khổng lồ), ép GCN phải định tuyến thông điệp dựa trên độ tin cậy thực tế.

            \item \textbf{Graph Isomorphism Network (GINE):} Nguyên bản GINE yêu cầu đặc trưng cạnh và nút phải đồng nhất về không gian. Để khắc phục sự bất đồng bộ giữa vector nút đa chiều $h_j$ và trọng số vô hướng $w_{ij}$, hệ thống ép kiểu tensor và thiết lập \texttt{edge\_dim = 1}. Thao tác này kích hoạt ma trận phóng chiếu nội tại $\mathbf{W}_e$, đẩy $w_{ij}$ vào cùng không gian ẩn và cộng trực tiếp vào láng giềng trước hàm kích hoạt phi tuyến:
            \[
                h_i^{(l+1)} = \text{MLP}^{(l)} \left( (1 + \epsilon) h_i^{(l)} + \sum_{j \in \mathcal{N}(i)} \text{ReLU} \left( h_j^{(l)} + \mathbf{W}_e w_{ij} \right) \right)
            \]
            Tinh chỉnh này đảm bảo mô hình không bị mất mát tín hiệu hành vi, kết hợp hoàn hảo giữa hình học tô-pô và định lượng hành vi trên mạng xã hội phi tập trung.

            \item \textbf{Principal Neighbourhood Aggregation (PNA):} Để khắc phục điểm mù về thông tin giao dịch của các toán tử thống kê, hệ thống áp dụng cơ chế dung hợp sớm thông qua mạng tiền xử lý \texttt{pre\_layers = 1} (kí hiệu là $\text{MLP}_{pre}$). "Công thức truyền tin được tùy chỉnh mở rộng:
            \[
                h_i^{(l+1)} = \text{MLP}_{post} \left( h_i^{(l)}, \bigoplus_{j \in \mathcal{N}(i)} \text{MLP}_{pre} \left( h_j^{(l)} \parallel w_{ij} \right) \right)
            \]
            Đặc trưng láng giềng $h_j$ buộc phải ghép nối và hòa trộn với $w_{ij}$ thành thông điệp hợp nhất trước khi đi qua khối thống kê $\bigoplus$. Nhờ vậy, độ lệch chuẩn ($\sigma$) sinh ra từ $\bigoplus$ chuyển sang đo lường sự biến thiên của lịch sử tương tác — một chỉ báo nhạy bén để bắt các mạng lưới Sybil có các hành vi rập khuôn.

            \item \textbf{Graph Transformer Networks:} Để đối phó với kịch bản tấn công sao chép - nơi bot có vector hồ sơ giống hệt người thật. Kế thừa cơ chế của mạng UniMP, hệ thống tiêm trực tiếp trọng số $w_{ij}$ vào cả ma trận Khóa (\textit{Key}) và Giá trị (\textit{Value}). Công thức tổng hợp láng giềng được cập nhật thành:
            \[
                h_i^{(l+1)} = \sum_{j \in \mathcal{N}(i)} \alpha_{ij} \left( \mathbf{W}_v h_j^{(l)} + \mathbf{W}_e w_{ij} \right)
            \]
            Trong đó, hệ số chú ý $\alpha_{ij}$ được điều chỉnh bởi chi phí giao dịch thông qua tích vô hướng:
            \[
                \alpha_{ij} = \text{softmax} \left( \frac{(\mathbf{W}_q h_i^{(l)})^T (\mathbf{W}_k h_j^{(l)} + \mathbf{W}_e w_{ij})}{\sqrt{d}} \right)
            \]
            Hệ quả là, dù vector hồ sơ $h_i$ và $h_j$ có độ trùng khớp cao đến đâu, các liên kết có chi phí thao tác rẻ tiền sẽ sinh ra giá trị $\mathbf{W}_e w_{ij}$ rất nhỏ, làm kéo tụt điểm tích vô hướng. Lõi Transformer sẽ tự động ngắt bỏ sự chú ý ($\alpha_{ij} \rightarrow 0$) đối với các nút láng giềng này.

            \end{itemize}
        \end{itemize}

Thông qua phần tối ưu hóa này, đề tài không chỉ giải quyết được bài toán về giới hạn phần cứng mà còn thiết lập một quy trình thống nhất.